% Full title: Chinese Translation of M\"uller Representation and the Coupled Operator
\section{M\"uller 表示与耦合算子}
\label{sec:muller-coupling}

\subsection{Rayleigh 展开与层势约定}

令 $\nu$ 为 $\Upsilon^0=\partial\Omega^0$ 上从 $\Omega^0$ 指向背景区域的单位法向量。内、外迹定义为
\[
  \gamma^{\mathrm i}u(z):=\lim_{h\downarrow0}u(z-h\nu(z)),
  \qquad
  \gamma^{\mathrm e}u(z):=\lim_{h\downarrow0}u(z+h\nu(z)).
\]
沿 $\nu$ 方向的法向导数也使用相同的上标。二维出射基本解为
\begin{equation}
  \Phi_\kappa(x,z):=\frac{\ii}{4}H_0^{(1)}(\kappa|x-z|).
  \label{eq:fundamental-solution}
\end{equation}
对 $x\notin\Upsilon^0$，定义
\begin{align}
  (S_\kappa\sigma)(x)&:=\int_{\Upsilon^0}\Phi_\kappa(x,z)\sigma(z)\dd s(z),
  \label{eq:single-layer}\\
  (D_\kappa\tau)(x)&:=\int_{\Upsilon^0}
  \partial_{\nu(z)}\Phi_\kappa(x,z)\tau(z)\dd s(z).
  \label{eq:double-layer}
\end{align}
在这些约定下，$\gamma_D^{\mathrm e}D_\kappa-\gamma_D^{\mathrm i}D_\kappa=I$，且 $\gamma_N^{\mathrm e}S_\kappa-\gamma_N^{\mathrm i}S_\kappa=-I$。

令 $G^{\mathrm{qp}}_{k,\beta}$ 为竖直方向 $\beta$-准周期的出射 Green 函数，满足
\[
  G^{\mathrm{qp}}_{k,\beta}(x+de_y,z)
  =e^{\ii\beta d}G^{\mathrm{qp}}_{k,\beta}(x,z),
  \qquad (\Delta_x+k^2)G^{\mathrm{qp}}_{k,\beta}(x,z)=-\delta_z.
\]
在其例外集合之外，它由格点和表示为
\begin{equation}
  G^{\mathrm{qp}}_{k,\beta}(x,z)
  =\sum_{\ell\in\Z}e^{\ii\beta\ell d}
  \Phi_k(x,z+\ell d e_y),
  \label{eq:qp-Green}
\end{equation}
必要时按标准的收敛延拓解释。把 \eqref{eq:single-layer}--\eqref{eq:double-layer} 中的 $\Phi_k$ 替换为 $G^{\mathrm{qp}}_{k,\beta}$，即可定义 $S^{\mathrm{qp}}_{k,\beta}$ 和 $D^{\mathrm{qp}}_{k,\beta}$。

令
\[
  \gamma_m(k,\beta):=\sqrt{k^2-\beta_m^2},
  \qquad \operatorname{Im}\gamma_m\geq0.
\]
若某个 $m$ 使 $\gamma_m=0$，则发生 Wood 异常。在非 Wood 参数处定义
\begin{equation}
  h^{1/2}_\beta:=
  \left\{a=(a_m)_{m\in\Z}:
  \sum_m\langle\beta_m\rangle|a_m|^2<\infty\right\},
  \qquad
  \Xi_\beta:=h^{1/2}_\beta\times h^{1/2}_\beta.
  \label{eq:Rayleigh-space}
\end{equation}
对 $\xi=(\xi^\rightarrow,\xi^\leftarrow)\in\Xi_\beta$，中心胞元内的齐次 Rayleigh 重构为
\begin{equation}
  (\mathcal H_{k,\beta}\xi)(x,y)
  :=\sum_{m\in\Z}\left[
  \xi_m^\rightarrow e^{\ii\gamma_m(x-X^-)}
  +\xi_m^\leftarrow e^{-\ii\gamma_m(x-X^+)}\right]\psi_m(y).
  \label{eq:Rayleigh-reconstruction}
\end{equation}
\eqref{eq:Rayleigh-space} 中的权重同时控制 $H^{1/2}_\beta$ Dirichlet 迹和 $H^{-1/2}_\beta$ 法向迹。

界面密度空间为
\begin{equation}
  \mathcal D^0:=H^{1/2}(\Upsilon^0)\times H^{-1/2}(\Upsilon^0),
  \qquad \eta=(\tau,\sigma)\in\mathcal D^0,
  \label{eq:density-space}
\end{equation}
其中 $\tau$ 是双层密度，$\sigma$ 是单层密度。为描述两个材料迹尚未匹配时的场，令
\begin{equation}
  H^1_{\beta,\mathrm{br}}(\mathcal C^0):=
  H^1_\beta(\mathcal C^0\setminus\overline{\Omega^0})
  \oplus H^1(\Omega^0).
  \label{eq:center-broken-space}
\end{equation}
下标 ``br'' 表示\emph{在 $\Upsilon^0$ 处断开}：在下述传输残差被置零之前，两个分量具有相互独立的迹。令 $k^0:=n^0k$。中心区重构
\begin{equation}
  \mathcal R_{\mathrm c}(k,\beta):\mathcal D^0\times\Xi_\beta
  \longrightarrow H^1_{\beta,\mathrm{br}}(\mathcal C^0)
  \label{eq:center-reconstruction-map}
\end{equation}
定义为断开场 $u_{\eta,\xi}^{\mathrm c}$：
\begin{equation}
  u_{\eta,\xi}^{\mathrm c}(x)=
  \begin{cases}
  \mathcal H_{k,\beta}\xi(x)
  +D^{\mathrm{qp}}_{k,\beta}\tau(x)
  -S^{\mathrm{qp}}_{k,\beta}\sigma(x),
  &x\in\mathcal C^0\setminus\overline{\Omega^0},\\
  D_{k^0}\tau(x)-S_{k^0}\sigma(x),
  &x\in\Omega^0.
  \end{cases}
  \label{eq:center-reconstruction}
\end{equation}
因此，外部场分量在竖直方向准周期，而有界内部场分量不需要单独的准周期条件。

所有迹均沿固定法向 $\nu$ 取值，由此定义 M\"uller 传输残差：
\begin{equation}
\begin{split}
  \mathcal M(k,\beta)(\eta,\xi):=
  \bigl(&\gamma_D^{\mathrm e}u_{\eta,\xi}^{\mathrm c}
         -\gamma_D^{\mathrm i}u_{\eta,\xi}^{\mathrm c},\\
        &\gamma_N^{\mathrm e}u_{\eta,\xi}^{\mathrm c}
         -\gamma_N^{\mathrm i}u_{\eta,\xi}^{\mathrm c}\bigr)
  \in\mathcal H_{\Upsilon^0},
  \qquad
  \mathcal H_{\Upsilon^0}:=H^{1/2}(\Upsilon^0)\times H^{-1/2}(\Upsilon^0).
\end{split}
\label{eq:Muller-residual}
\end{equation}
这一迹定义无需依赖未写明的矩阵约定，就固定了每个半单位项的符号。利用跳跃关系展开它，即得到 M\"uller 块。令
\begin{equation}
  \mathcal Z_{\mathrm c}(k,\beta):=\ker\mathcal M(k,\beta),
  \qquad
  \mathcal N_{\mathrm c}(k,\beta):=
  \mathcal Z_{\mathrm c}(k,\beta)\cap\ker\mathcal R_{\mathrm c}(k,\beta).
  \label{eq:center-nullspaces}
\end{equation}
因此，$\mathcal N_{\mathrm c}$ 恰好包含所有满足传输方程、却生成零物理中心区场的中心区代数表示。

令 $\mathcal S_{\mathrm c}(k,\beta)$ 为 $H^1_\beta(\mathcal C^0)$ 中满足 $\Upsilon^0$ 上物理 TM 传输条件、但在 $\Gamma^\pm$ 上不施加边界条件的分片 Helmholtz 场空间。定义表示例外集合 $\mathfrak E_{\mathrm{rep}}(\beta)$ 为以下集合之并：
\begin{enumerate}[label=(\roman*),leftmargin=2.2em]
  \item Wood 参数 $\gamma_m(k,\beta)=0$；
  \item 所选准周期 Green 函数的极点或未定义参数；
  \item 互补交换波数传输问题不唯一可解的参数，其中内部波数为 $k$、外部波数为 $k^0$。
\end{enumerate}
即便参数不在这一明确列出的集合中，下文商空间仍保留可能存在的表示零空间。

\begin{theorem}[中心胞元 M\"uller--Rayleigh 表示]
\label{thm:center-representation}
设 $k^2\in I_\beta$ 且 $k\notin\mathfrak E_{\mathrm{rep}}(\beta)$。$\mathcal R_{\mathrm c}(k,\beta)$ 在 $\mathcal Z_{\mathrm c}(k,\beta)$ 上的限制满射到 $\mathcal S_{\mathrm c}(k,\beta)$。因此，它诱导线性同构
\begin{equation}
  \mathcal Z_{\mathrm c}(k,\beta)/\mathcal N_{\mathrm c}(k,\beta)
  \xrightarrow{\ \cong\ }\mathcal S_{\mathrm c}(k,\beta),
  \qquad [\eta,\xi]\longmapsto\mathcal R_{\mathrm c}(k,\beta)(\eta,\xi).
  \label{eq:center-quotient-isomorphism}
\end{equation}
这里不主张密度唯一。密度唯一性等价于附加结论 $\mathcal N_{\mathrm c}(k,\beta)=\{0\}$。
\end{theorem}
\status{适配与证明目标。}
\begin{proofroadmap}
在 $\Gamma^-$ 和 $\Gamma^+$ 的齐次邻域中，从给定中心区场提取两个齐次 Rayleigh 分量。减去它们按 \eqref{eq:Rayleigh-reconstruction} 所作的延拓；余项满足出射单胞元表示的假设。应用准周期外部 Green 恒等式和自由空间内部 Green 恒等式，并依照 \cite[Theorem~4 and Appendix~A, Lemmas~9--10]{BarnettGreengard2010} 构造互补交换波数场。在该计算中，外部双层算子是 $D^{\mathrm{qp}}_{k,\beta}$，夹杂双层算子则是 $D_{k^0}$；在外部消去恒等式中使用 $D^{\mathrm{qp}}_{k^0,\beta}$ 会造成核不匹配，必须避免。

物理场映射的满射性是实质性结论。一旦证明它，商空间同构便由第一同构定理推出。密度核必须借助跳跃关系、互补场和 Calder\'on 投影范围另行分析；\cite[Lemmas~2.2--2.6]{HiptmairMoiolaSpence2022} 中的传输边界积分方程分析说明，物理解的唯一性定理并不能自动推出密度唯一性。

详细内部指引见 \path{research/planning/muller-cauchy/notes/layers.md}，以及 \path{research/planning/muller-cauchy/p-proofs.md} 中名为 ``corrected center representation'' 和 ``density-to-field kernel characterization'' 的待证事项。最低可接受结果是模去一个有限维遗漏的齐次扇区后得到满射；只有在所述正则集合上证明 $\mathcal N_{\mathrm c}=\{0\}$ 后，才能接受不取商的定理。
\end{proofroadmap}
\begin{proof}
写
$u=(u^{\mathrm e},u^{\mathrm i})\in\mathcal S_{\mathrm c}(k,\beta)$，
其中两项分别是外部与夹杂内部的分量。我们构造
$(\eta,\xi)\in\mathcal Z_{\mathrm c}(k,\beta)$，使其重构场恰为 $u$。

\emph{步骤 1：分离入射的齐次 Rayleigh 场。}
由于 $\Omega^0\Subset\mathcal C^0$，两个端口都具有齐次材料邻域。在任一
这样的邻域中，$u^{\mathrm e}$ 关于模态 $\psi_m$ 的 Fourier 系数满足
\[
  \partial_x^2u_m+\gamma_m^2u_m=0.
\]
非 Wood 假设给出 $\gamma_m\neq0$，故该解具有唯一的右行/左行分解。记左端口
邻域中的右行系数为 $a_m^{-,\rightarrow}$，右端口邻域中的左行系数为
$a_m^{+,\leftarrow}$；相位基点分别取为 $X^-$ 和 $X^+$。若
$(f_m^\pm,p_m^\pm)$ 分别是两个端口上 Dirichlet 迹与 $\partial_x$ 迹的
Fourier 系数，则
\[
  a_m^{-,\rightarrow}
  =\frac12\left(f_m^-+\frac{p_m^-}{\ii\gamma_m}\right),
  \qquad
  a_m^{+,\leftarrow}
  =\frac12\left(f_m^+-\frac{p_m^+}{\ii\gamma_m}\right).
\]
迹定理给出 $f^\pm\in h^{1/2}_\beta$ 与
$p^\pm\in h^{-1/2}_\beta$。此外，对充分大的 $|m|$ 有
$|\gamma_m|\asymp\langle\beta_m\rangle$；剩余指标只有有限多个，并且其中
没有 $\gamma_m=0$。因此
\[
  \xi:=\bigl((a_m^{-,\rightarrow})_m,
             (a_m^{+,\leftarrow})_m\bigr)\in\Xi_\beta.
\]
令 $h:=\mathcal H_{k,\beta}\xi$，$w:=u^{\mathrm e}-h$。在左端口邻域中，
$w$ 只含左行模态；在右端口邻域中，$w$ 只含右行模态，因为从另一端口发出的
$h$ 分量在本端口只贡献相应的出射分量。将这些模态级数穿过两个端口继续延拓，
便得到 $w$ 在 $B\setminus\overline{\Omega^0}$ 上的延拓；它满足竖直方向
$\beta$-准周期的方程 $(\Delta+k^2)w=0$，并在两个无穷端均为出射场。

\emph{步骤 2：两组 Green 表示恒等式。}
始终使用固定法向量 $\nu$，并简记
\[
  P_\kappa(f,q):=D_\kappa f-S_\kappa q,
  \qquad
  P^{\mathrm{qp}}_{k,\beta}(f,q)
  :=D^{\mathrm{qp}}_{k,\beta}f-S^{\mathrm{qp}}_{k,\beta}q.
\]
先在截断的准周期条带上应用 Green 第二恒等式，再通过极限吸收过程取出射极限，
得到
\[
  P^{\mathrm{qp}}_{k,\beta}
  (\gamma_D^{\mathrm e}w,\gamma_N^{\mathrm e}w)
  =w
  \quad\text{于 }B\setminus\overline{\Omega^0}.
\]
对 $\Omega^0$ 中满足 $(\Delta+k^2)z=0$ 的解 $z$ 应用同一恒等式，得到消去关系
\[
  P^{\mathrm{qp}}_{k,\beta}
  (\gamma_D^{\mathrm i}z,\gamma_N^{\mathrm i}z)=0
  \quad\text{于 }B\setminus\overline{\Omega^0}.
\]
对于波数 $k^0$ 的自由空间核，相应恒等式为
\begin{align*}
  P_{k^0}(\gamma_D^{\mathrm i}z,\gamma_N^{\mathrm i}z)
  &=-z&&\text{于 }\Omega^0,\\
  P_{k^0}(\gamma_D^{\mathrm e}z,\gamma_N^{\mathrm e}z)
  &=0&&\text{于 }\Omega^0,
\end{align*}
其中后一式中的 $z$ 是出射外部场。符号由 \eqref{eq:single-layer} 前规定的跳跃
约定确定。这些恒等式是
\cite[Appendix~A, Lemmas~9--10]{BarnettGreengard2010} 所用恒等式的单周期
版本：竖直边界项由 $\beta$-准周期性抵消，两个无穷端的项则由出射极限吸收
过程消去。

\emph{步骤 3：构造公共 M\"uller 密度。}
在 $\Upsilon^0$ 上令
\[
  f:=\gamma_D^{\mathrm e}u^{\mathrm e}
    =\gamma_D^{\mathrm i}u^{\mathrm i},
  \qquad
  q:=\gamma_N^{\mathrm e}u^{\mathrm e}
    =\gamma_N^{\mathrm i}u^{\mathrm i},
\]
并记 $h_D:=\gamma_Dh$、$h_N:=\gamma_Nh$。于是 $w$ 的外部 Cauchy 数据为
$(f-h_D,q-h_N)$。考虑下列互补交换波数传输问题：
\begin{align*}
  (\Delta+k^2)v^{\mathrm i}&=0
    &&\text{于 }\Omega^0,\\
  (\Delta+(k^0)^2)v^{\mathrm e}&=0
    &&\text{于 }\R^2\setminus\overline{\Omega^0},
\end{align*}
其中 $v^{\mathrm e}$ 满足 Sommerfeld 条件，并具有跳跃
\begin{align*}
  \gamma_D^{\mathrm e}v^{\mathrm e}
  -\gamma_D^{\mathrm i}v^{\mathrm i}&=2f-h_D,\\
  \gamma_N^{\mathrm e}v^{\mathrm e}
  -\gamma_N^{\mathrm i}v^{\mathrm i}&=2q-h_N.
\end{align*}
由 $\mathfrak E_{\mathrm{rep}}(\beta)$ 定义中的第 (iii) 项以及
$k\notin\mathfrak E_{\mathrm{rep}}(\beta)$，该问题对这些允许的迹数据具有唯一解。
定义
\begin{align*}
  \tau
  &:=f-h_D+\gamma_D^{\mathrm i}v^{\mathrm i}
    =-f+\gamma_D^{\mathrm e}v^{\mathrm e},\\
  \sigma
  &:=q-h_N+\gamma_N^{\mathrm i}v^{\mathrm i}
    =-q+\gamma_N^{\mathrm e}v^{\mathrm e}.
\end{align*}
这两组等式恰由所指定的两个跳跃得到；迹映射性质还给出
$(\tau,\sigma)\in\mathcal D^0$。

令 $\eta=(\tau,\sigma)$。利用上面的两个准周期恒等式，重构场的外部分量为
\begin{align*}
  h+P^{\mathrm{qp}}_{k,\beta}(\tau,\sigma)
  &=h+P^{\mathrm{qp}}_{k,\beta}(f-h_D,q-h_N)\\
  &\quad
    +P^{\mathrm{qp}}_{k,\beta}
      (\gamma_D^{\mathrm i}v^{\mathrm i},
       \gamma_N^{\mathrm i}v^{\mathrm i})\\
  &=h+w=u^{\mathrm e}.
\end{align*}
再利用 $\tau,\sigma$ 的另一组表达式和上面的两个自由空间恒等式，重构场的
内部分量为
\begin{align*}
  P_{k^0}(\tau,\sigma)
  &=-P_{k^0}(f,q)
    +P_{k^0}(\gamma_D^{\mathrm e}v^{\mathrm e},
             \gamma_N^{\mathrm e}v^{\mathrm e})\\
  &=u^{\mathrm i}.
\end{align*}
因此 $\mathcal R_{\mathrm c}(k,\beta)(\eta,\xi)=u$。由于 $u$ 满足物理传输条件，
它的两个跳跃均为零，故
$\mathcal M(k,\beta)(\eta,\xi)=0$，且
$(\eta,\xi)\in\mathcal Z_{\mathrm c}(k,\beta)$。满射性得证。

\emph{步骤 4：转到商空间。}
由 \eqref{eq:center-nullspaces}，
$\mathcal R_{\mathrm c}|_{\mathcal Z_{\mathrm c}}$ 的核恰为
$\mathcal N_{\mathrm c}$。第一同构定理遂给出
\eqref{eq:center-quotient-isomorphism}。这一论证既不需要也不主张
$\mathcal N_{\mathrm c}$ 为零。
\end{proof}

\subsection{连续耦合算子}

中心区重构的端口迹是有界映射
\begin{equation}
  \Pi^\pm(k,\beta):=\Tr^{0,\pm}\mathcal R_{\mathrm c}(k,\beta):
  \mathcal D^0\times\Xi_\beta\longrightarrow\mathcal H(\Gamma^\pm).
  \label{eq:center-port-map}
\end{equation}
该定义同时包含层势贡献与齐次 Rayleigh 贡献。半波导映射 $\mathcal F^\pm$ 和 $Q^\pm$ 已在 \eqref{eq:field-synthesis}--\eqref{eq:trace-synthesis} 中定义。

\begin{definition}[耦合 M\"uller--Cauchy 算子与表示零空间]
\label{def:coupled-operator}
定义代数定义域与残差值域为
\begin{align}
  \mathcal X(k,\beta)
  &:=\mathcal D^0\times\Xi_\beta\times
  \mathcal K_s^-\times\mathcal K_s^+,
  \label{eq:coupled-domain}\\
  \mathcal Y(k,\beta)
  &:=\mathcal H_{\Upsilon^0}\times
  \mathcal H(\Gamma^-)\times\mathcal H(\Gamma^+).
  \label{eq:coupled-codomain}
\end{align}
对 $x=(\eta,\xi,c^-,c^+)\in\mathcal X(k,\beta)$，令
\begin{equation}
  \mathbb A(k,\beta)x:=
  \begin{bmatrix}
  \mathcal M(k,\beta)(\eta,\xi)\\
  \Pi^-(k,\beta)(\eta,\xi)-Q^-(k,\beta)c^-\\
  \Pi^+(k,\beta)(\eta,\xi)-Q^+(k,\beta)c^+
  \end{bmatrix}
  \in\mathcal Y(k,\beta).
  \label{eq:coupled-operator}
\end{equation}
第一行是 $\Upsilon^0$ 上的物理传输失配；第二、三行分别是左、右端口上完整的面向中心区 Cauchy 数据失配。

令
\begin{equation}
  H^1_{\beta,\mathrm{br}}(B):=
  H^1_\beta(W^- )\oplus
  H^1_\beta(\mathcal C^0\setminus\overline{\Omega^0})
  \oplus H^1(\Omega^0)
  \oplus H^1_\beta(W^+).
  \label{eq:global-broken-space}
\end{equation}
全局断开场重构为
\begin{equation}
  \mathcal R_{\mathrm g}(k,\beta)x:=
  \begin{cases}
    \mathcal F^-(k,\beta)c^-,&\text{于 }W^-,\\
    \mathcal R_{\mathrm c}(k,\beta)(\eta,\xi),&\text{于 }\mathcal C^0,\\
    \mathcal F^+(k,\beta)c^+,&\text{于 }W^+.
  \end{cases}
  \label{eq:global-reconstruction}
\end{equation}
对一般 $x$，它是断开场；当 $x\in\ker\mathbb A(k,\beta)$ 时，它是全局 $H^1_\beta(B)$ 场。最后，定义表示诱发的零空间
\begin{equation}
  \mathcal N_{\mathrm{rep}}(k,\beta):=
  \left\{x\in\ker\mathbb A(k,\beta):
  \mathcal R_{\mathrm g}(k,\beta)x=0
  \text{ 于 }H^1_{\beta,\mathrm{br}}(B)\right\}.
  \label{eq:Nrep-definition}
\end{equation}
因此，$\mathbb A$ 是完全定义的代数耦合，而 $\mathcal N_{\mathrm{rep}}$ 精确记录其核中不承载任何物理场的向量。这两个符号均不是某个未明确指定的数值矩阵的简写。
\end{definition}
